\section*{Acknowledgments}
\addcontentsline{toc}{section}{Acknowledgments}

Foremost, I am deeply grateful to my thesis advisor, Prof.\ Bj\"orn Sandstede,
who shaped this project at every level.
From our earliest conversations about what question was even worth asking,
through countless meetings where he pointed me toward the right framework,
to the patient feedback that sharpened every chapter, his guidance has been
indispensable.
His instinct for identifying the core mathematical structure behind a messy
numerical experiment --- and his trust that I could work my way toward it ---
gave me both the direction and the confidence to see this work through.
I am also grateful to my second reader, Prof.\ Ian~Wong, for his critical
perspective and encouragement.

This thesis is, in many ways, a synthesis of everything I studied at Brown
in the APMA--CS--Mathematics track.
The pipeline draws on probability theory and statistics (KDE density fields,
noise models, convergence analysis), linear algebra and functional analysis
(POD, SVD truncation, Hilbert-space weak formulations), differential equations
(ETDRK4 time-stepping, PDE library construction, the Fokker--Planck framing),
machine learning and deep learning (LSTM sequence models, SR3 sparse
regression), and even abstract algebra --- group theory enters through the
continuous SO(2) rotational symmetry that motivates the shift-alignment step,
where symmetry reduction has the flavour of quotienting by a group action.
Topology threads through the persistent homology literature we position
against, and the Kolmogorov $n$-width argument that motivates nonlinear model
reduction is a piece of functional analysis.
I am grateful to all the faculty and courses at Brown that gave me the tools
to build something at this intersection.

I thank the Center for Computation and Visualization
(CCV) at Brown University for providing access to
the OSCAR high-performance computing
cluster, on which all particle simulations,
ROM training, WSINDy fitting, and ETDRK4 rollout computations were
performed.
